

% choose options for [] as required from the list
% in the Reference Guide

%\usepackage{mathptmx}
%\usepackage{helvet}
%\usepackage{courier}
%


\usepackage{type1cm}

\usepackage{makeidx}         % allows index generation
\usepackage{graphicx}        % standard LaTeX graphics tool
                             % when including figure files
\usepackage{multicol}        % used for the two-column index
\usepackage{needspace}       % keep selected statements together after page review
\usepackage[bottom]{footmisc}% places footnotes at page bottom

\usepackage{newtxtext}       %
\usepackage{newtxmath}       % selects Times Roman as basic font

\usepackage[unicode]{hyperref}
\usepackage[nameinlink,capitalize]{cleveref}
\crefname{appendix}{Appendix}{Appendices}
\Crefname{appendix}{Appendix}{Appendices}
\hypersetup{
colorlinks=true,
linktoc=all
}
\makeatletter
% Share the section theorem counter while retaining each environment's reference type.
\@for\bookenv:=conjecture,corollary,definition,example,lemma,problem,proposition,question,remark\do{%
  \expandafter\let\csname c@\bookenv\endcsname\c@theorem
}
\renewcommand*{\theHtheorem}{\theHsection.\arabic{theorem}}
\makeatother

\usepackage{enumitem}
\usepackage{mathtools}
\providecommand{\xlongrightarrow}[2][]{\xrightarrow[#1]{#2}}
\usepackage{tikz-cd}
\usetikzlibrary{decorations.pathreplacing,arrows.meta,calc,patterns}
\usepackage{dsfont}

\usepackage{fancyhdr}
\renewcommand{\headrulewidth}{0pt}

\usepackage[final]{microtype}
\usepackage{inputenc}
\usepackage{textcomp}
\usepackage[T1]{fontenc}

%\usepackage[color]{showkeys}
%\definecolor{refkey}{rgb}{0,0,1}


% see the list of further useful packages
% in the Reference Guide

\AddToHook{env/theindex/begin}{\raggedcolumns\setlength{\maxbalancingoverflow}{0pt}}
\makeindex             % used for the subject index
                       % please use the style svind.ist with
                       % your makeindex program

%%%%%%%%%%%%%%%%%%%%%%%%%%%%%%%%%%%%%%%%%%%%%%%%%%%%%%%%%%%%%%%%%%%%%

\graphicspath{{./}{images/}}

\newcommand{\nef}{\mathrm{nef}}
\newcommand{\Bir}{\mathrm{Bir}}
\newcommand{\NS}{\mathrm{NS}}
\newcommand{\lex}{\mathrm{lex}}
\newcommand{\NA}{\mathrm{NA}}
\newcommand{\SH}{\mathrm{SH}}
\newcommand{\PSH}{\mathrm{PSH}}
\newcommand{\QPSH}{\mathrm{QPSH}}
\newcommand{\BM}{\mathrm{BM}}
\newcommand{\Conv}{\mathrm{Conv}}
\newcommand{\FS}{\mathrm{FS}}
\newcommand{\An}{\mathrm{An}}
\newcommand{\val}{\mathrm{val}}
\newcommand{\Macron}{\mathrm{Defect}}
\newcommand{\Bl}{\mathrm{Bl}}
\newcommand{\Ele}{\mathrm{Ele}}
\newcommand{\triv}{\mathrm{triv}}
\newcommand{\gf}{\mathrm{gf}}
\newcommand{\tor}{\mathrm{tor}}
\newcommand{\prop}{\mathrm{prop}}
\newcommand{\Star}{\mathrm{Star}}
\newcommand{\Haus}{\mathrm{Haus}}
\newcommand{\sg}{\mathrm{sg}}
\newcommand{\TC}{\mathrm{TC}}
\newcommand{\Kah}{\textup{Käh}}
\newcommand{\sups}{\operatorname*{sup*}}
\newcommand{\uscsup}[1][]{\operatorname*{sup}_{#1}{}\!{}^{*}}
\newcommand{\ddc}{\mathrm{dd}^{\mathrm{c}}}
\newcommand{\ddt}{\frac{\mathrm{d}}{\mathrm{d}t}}
\newcommand{\ddtz}[1][0]{\left.\ddt\right|_{t={#1}}}
\newcommand{\loc}{\mathrm{loc}}
\newcommand{\convC}{\xrightarrow{\mathrm{C}}}
\newcommand{\qsup}{\operatorname*{q-sup}}
\newcommand{\qinf}{\operatorname*{q-inf}}

\DeclareMathOperator{\Img}{Im}
\DeclareMathOperator{\Nm}{Norm}
\DeclareMathOperator{\nn}{nn}
\DeclareMathOperator{\nK}{nK}
\DeclareMathOperator{\Capa}{Cap}
\DeclareMathOperator{\Real}{Re}
\DeclareMathOperator{\Spec}{Spec}
\DeclareMathOperator{\Sing}{Sing}
\DeclareMathOperator{\Sgn}{Sign}
\DeclareMathOperator{\Rest}{Res}
\DeclareMathOperator{\DHm}{DH}
\DeclareMathOperator{\Fin}{Fin}
\DeclareMathOperator{\Reg}{Reg}
\DeclareMathOperator{\Tr}{Tr}
\DeclareMathOperator{\Trop}{Trop}
\DeclareMathOperator{\Div}{div}
\DeclareMathOperator{\ord}{ord}
\DeclareMathOperator{\Dom}{Dom}
\DeclareMathOperator{\epi}{epi}
\DeclareMathOperator{\LCE}{CE}
\DeclareMathOperator{\Supp}{Supp}
\DeclareMathOperator{\MA}{MA}
\DeclareMathOperator{\vol}{vol}
\DeclareMathOperator{\Pic}{Pic}
\DeclareMathOperator{\Int}{Int}
\DeclareMathOperator{\RelInt}{RelInt}
\DeclareMathOperator{\cl}{cl}
\DeclareMathOperator{\cor}{cor}
\DeclareMathOperator{\GL}{GL}
\DeclareMathOperator{\Aut}{Aut}
\DeclareMathOperator{\usc}{usc}
\DeclareMathOperator{\mult}{mult}
\DeclareMathOperator{\Bs}{Bs}




\setlist[enumerate,1]{label=\textup{(\arabic*)},ref=\textup{(\arabic*)}}

\makeatletter
% For labels of the form itm:<parent label>:<item>, print "Parent (item)".
\newcommand{\itemref}[1]{\itemref@split\cref#1:\@nil}
\newcommand{\Itemref}[1]{\itemref@split\Cref#1:\@nil}
\def\itemref@split#1itm:#2:#3:#4:\@nil{#1{#2:#3}\,\ref{itm:#2:#3:#4}}
\makeatother



% --- Hyphenation hints for technical terms (added during indexing pass) ---
\hyphenation{
  plu-ri-sub-har-mon-ic plu-ri-sub-har-mon-ic-i-ty plu-ri-po-lar plu-ri-po-ten-tial
  pseu-do-con-vex pseu-do-met-ric pseu-do-norm sub-har-mon-ic
  bi-mero-mor-phic mero-mor-phic ho-lo-mor-phic
  semi-con-tin-u-ous semi-con-ti-nu-i-ty up-per-semi-con-ti-nu-i-ty
  multi-plier multi-pli-er ana-lyt-i-fi-ca-tion
  equi-sin-gu-lar equi-sin-gu-lar-i-ty quasi-equi-sin-gu-lar
  trop-i-cal-i-za-tion non-Arch-i-me-de-an Arch-i-me-de-an
  sub-ge-o-des-ic sub-ge-o-des-ics ge-o-des-ic ge-o-des-ics
  Mon-ge--Am-pe-re Käh-ler Lel-ong Bo-uck-som Ohsa-wa--Take-go-shi
  Da-r-vas Bed-ford--Tay-lor Don-ald-son Ross--Witt-Nys-tröm
  Bouck-som--Eys-sid-ieux--Guedj--Zer-i-ahi
  ho-mo-ge-ne-ous ho-mo-ge-ne-i-ty Berk-o-vich
  Brel-ot-Grau-ert--Rem-mert Hi-ron-a-ka Wein-stock
  Lazars-feld Ka-veh--Khov-an-skii Kis-el-man Choq-uet
}
% --- end hyphenation hints ---

%\includeonly{author/chapter_pshfunctions,author/chapter_npp,author/chapter_envelopes,author/chapter_Bergman}
%\includeonly{author/chapter_envelopes,author/chapter_comp,author/chapter_Bergman}

